%%%%%%%%%%  scan idx 213  |  folio 200  %%%%%%%%%%
\[
  \Gamma(U'\times V',{}^r\cO_{U'\times V'})=C(U',C^r(V')),
\]
où $U'$ (resp. $V'$) est un ouvert de $U$ (resp. $V$), et $C(U',C^r(V'))$ désigne l'espace des applications continues de $U'$ dans $C^r(V')$. On a une injection canonique
\[
  \pr_1^{-1}(\cO_U)\longrightarrow{}^r\cO_{U\times V},
\]
qui fait de $(U\times V,{}^r\cO_{U\times V})$ un espace annelé au-dessus de $U$.

On dit qu'un espace annelé $X$ au-dessus de $S$ est une \emph{variété mixte de classe} $C^r$ ($r\in\N\cup\{\infty\}$) si $X$ est localement isomorphe (en tant qu'espace annelé au-dessus de $S$) à un espace du type $(U\times V,{}^r\cO_{U\times V})$. (Si $S$ est réduit à un point, une variété mixte de classe $C^r$ au-dessus de $S$ n'est autre qu'une variété de classe $C^r$ au sens usuel).

Soit $f:X\to S$ une variété mixte de classe $C^r$. Pour $s\in S$, la fibre $X_s=f^{-1}(s)$, munie du faisceau $i_s^{-1}(\cO_X)/I_s$, où $i_s:X_s\to X$ est l'injection et $I_s$ l'Idéal des sections qui «~s'annulent~» sur $X_s$, est une variété de classe $C^r$.

\subsubsection*{1.2. Opérateurs différentiels «~relatifs~»}
Il est facile de faire, sur le modèle de la Géométrie Algébrique (EGA IV 16, SGA 3 VII), une étude infinitésimale des variétés mixtes. Nous nous bornerons à passer rapidement en revue les notions dont nous aurons besoin.

\textbf{1.2.1.} Soit $f:X\to S$ une variété mixte de classe $C^\infty$ (1.1). Pour tout $n\in\N$, on définit le faisceau $P^n_{X/S}$ des \emph{parties principales à l'ordre} $n$ sur $X$. C'est de deux façons différentes, correspondant aux deux projections $\pr_1$ et $\pr_2$ de $X\times_S X$ sur $X$, une $\cO_X$-Algèbre augmentée vers $\cO_X$, localement libre de type fini en tant que $\cO_X$-Module. Quand nous ne préciserons pas la structure de $\cO_X$-Module considérée sur $P^n_{X/S}$, il sera sous-entendu qu'il s'agit de celle provenant de $\pr_1$. On a des isomorphismes canoniques
